% !TEX root = main.tex
% LOG (2026-10-27): the "H-based projection" extension, dropped from contributions.tex (Sec. 4.4).
% Reason: it removes the dilation mode of H by a linear constraint and returns a site-independent
% radial projection onto C, but it still fails (merged shifted sites, no power diagram) when the
% projected targets fall in a flat face of C, so it does not solve the problem that H_alpha addresses.
% Not compiled; kept for reference only. Labels below are no longer defined in the paper.

%% ---- Paragraph removed from Sec. 4.4 (Extensions) ----
\paragraph{$\jointdual$-based projection.}
\label{sec:gauge-fixed}
As discussed in \S\ref{sec:regularized-joint}, maximizing $\jointdual_\alpha$ adjusts the prescribed centroids $\tcentroid$ depending on the choice of the sites $\qq$.
We now introduce a modified formulation that projects $\tcentroid$ onto $\achievable$ independently of the sites.
Recall from \S\ref{sec:joint-constraints} that the maximizers of $\jointdual$, when they exist, are only defined up to the homotheties of the shifted sites, which $\jointdual_\alpha$ removes by penalizing the shifts. % at the cost of a dependence on the sites.
As shown in Appendix~\ref{sec:appendix-existence}, only the dilation mode affects $\jointdual$ when the target centroids are not realizable, since the translations and the additive constant leave $\jointdual$ unchanged under the identities in \eqref{eq:partition-feasibility}.
This motivates a construction that fixes the scale of the shifted sites directly, much like a gauge fixing, by maximizing $\jointdual$ subject to a single linear constraint on the shifts:
\begin{equation}
  \max_{\shift, \weight} \ \jointdual(\shift, \weight)
  \quad \text{s.t.} \quad
  \sum\nolimits_i \pvol_i\,\shift_i\tr (\tcentroid_i - \domainCentroid) = 0 ,
  \label{eq:gauge-fixed}
\end{equation}
where the targets $\tcentroid$ are assumed to satisfy \eqref{eq:partition-feasibility-centroid} and to not all coincide with $\domainCentroid$.
This constraint pins the overall scale of the shifted sites, so it rules out the dilation mode while leaving the translations and the additive constant. %, which are harmless gauge freedoms since they do not change the cells.
% Numerically, Newton's method solves \eqref{eq:gauge-fixed} with the sparse facet-based Hessian of \S\ref{sec:joint-constraints}, bordered by a single row and column with entries $\smash{2\,\pvol_i\,(\tcentroid_i - \domainCentroid)}$ on the shifts.
% 
To characterize the solution, we introduce a scalar multiplier $r \!\in\! \R$ for the linear constraint in \eqref{eq:gauge-fixed}.
This constraint only modifies the shift gradient \eqref{eq:joint-gradient-shift}, by the term $\smash{2\,r\,\pvol_i (\tcentroid_i - \domainCentroid)}$, so, with exact volumes, the stationary centroids become
\begin{equation}
  \centroid_i = \domainCentroid + (1 + r)(\tcentroid_i - \domainCentroid) ,
  \label{eq:gauge-centroids}
\end{equation}
\ie the targets uniformly dilated about the domain centroid.
Since $\achievable$ is convex and compact, and contains $\domainCentroid\mathbf{1}$, reached at $r \!=\! -1$, in its relative interior \cite[Thm.~3.6]{Bourne2026ConvexOrder}, the dilations for which \eqref{eq:gauge-centroids} remains achievable form a closed interval around $r \!=\! -1$.
At the solution, the objective is the transport cost from the sites to a partition with volumes $\pvol$, which depends on the cells only through their centroids (Appendix~\ref{sec:appendix-regularized}).
Substituting \eqref{eq:gauge-centroids} into this cost reduces \eqref{eq:gauge-fixed}, up to a constant, to
\begin{equation}
  \min_{r \in \R} \ - r \sum\nolimits_i \pvol_i\,\qq_i\tr (\tcentroid_i - \domainCentroid)
  \quad \text{s.t.} \quad
  \centroid \in \achievable .
  \label{eq:gauge-solution}
\end{equation}
This reveals the role of the sites, whose influence is reduced to the sign of a single scalar: since \eqref{eq:gauge-solution} is linear in $r$, a positive sum selects the largest feasible dilation, which preserves the configuration of the targets, while a negative sum selects the smallest one, which reflects it through $\domainCentroid$, and a zero sum leaves $r$ undetermined.
We therefore choose sites with a positive sum, such as $\qq = \tcentroid$, for which the sum equals $\smash{\sum_i \pvol_i \norm{\tcentroid_i - \domainCentroid}^2 > 0}$ by \eqref{eq:partition-feasibility-centroid}.
Beyond this sign, the returned centroids \eqref{eq:gauge-centroids} do not depend on the sites.

This construction thus defines a projection map: it radially projects the targets about $\domainCentroid$ onto the relative boundary of $\achievable$, dilating targets inside $\achievable$ ($r \!>\! 0$), contracting targets outside $\achievable$ ($r \!<\! 0$), and returning realizable targets exactly ($r \!=\! 0$).
As a limitation, the relative boundary of $\achievable$ also contains flat faces exposed by coinciding generators (\S\ref{sec:background}): when the projected targets fall in the relative interior of such a face, the maximizer has merged shifted sites, and the optimal partition is no longer a power diagram.

%% ---- Sentence removed from Sec. 4.5, "Link to Bourne2025Inverting" (end of paragraph) ----
The $\jointdual$-based projection of \S\ref{sec:extensions} instead removes the dilation mode with a single linear constraint, which keeps a convex problem for all targets, and returns the radially dilated targets \eqref{eq:gauge-centroids} rather than the nearest ones.

%% ---- Sentences removed from Sec. 4.5, "Link to Bourne2026ConvexOrder" (before the table pointer) ----
For targets in $\achievable$, as $\lambda$ decreases to the largest dilation factor $1 + r$ of \eqref{eq:gauge-solution}, their projection tends to the dilated targets \eqref{eq:gauge-centroids} while their generators shrink to zero.
The $\jointdual$-based projection of \S\ref{sec:extensions} computes this limit directly, without vanishing generators, and recovers their exact limit $\lambda \!\to\! 1$ for realizable targets.

%% ---- Row removed from Table tab:gauge-comparison ----
    $\jointdual$ on a hyperplane (\S\ref{sec:extensions}) & radial dilation \eqref{eq:gauge-centroids} & yes & always \\
